\normaltrue \difficilefalse \tdifficilefalse
\correctiontrue

\exer{Pompe oscillante $\star$ \label{GEO:03:C2:06:13}}
\setcounter{question}{0}\marginnote{\xpComp{GEO}{03}}
\index{Compétence C2-06}\index{Compétence GEO-03}
\index{Pompe oscillante}

\ifcorrection
\else
\marginnote{\textbf{Pas de corrigé pour cet exercice.}}
\fi

\ifprof
\else
Soit le mécanisme suivant. On a $\vect{AB}=R\vect{i_1}$ et $\vect{CA}=H\vect{j_0}$.
De plus, $R=\SI{40}{mm}$ et $H=\SI{60}{mm}$.
Par ailleurs, on note $\vect{CB}=\lambda(t)\vect{i_2}$.

\begin{marginfigure}
\includegraphics[width=\linewidth]{13_01}
\end{marginfigure}
\fi

\question{Tracer le graphe des liaisons.}

\ifprof
\begin{marginfigure}
\includegraphics[width=.4\linewidth]{13_01_c}
\end{marginfigure}
\fi

\question{Exprimer $\lambda(t)$ en fonction de $\theta(t)$.}

\ifprof
En réalisant une fermeture géométrique, on a
\[
\vect{AB}+\vect{BC}+\vect{CA}=\vect{0},
\]
soit
\[
R\vect{i_1}-\lambda(t)\vect{i_2}+H\vect{j_0}=\vect{0}.
\]
La projection dans le repère $\rep{0}$ donne
\[
\left\{
\begin{array}{l}
R\cos\theta(t)-\lambda(t)\cos\varphi(t)=0,\\
R\sin\theta(t)-\lambda(t)\sin\varphi(t)+H=0.
\end{array}
\right.
\]
Après élévation au carré et sommation, on obtient
\[
\lambda^2(t)=R^2+H^2+2HR\sin\theta(t).
\]
La longueur étant positive,
\[
\boxed{\lambda(t)=\sqrt{R^2+H^2+2HR\sin\theta(t)}}.
\]
\fi

\question{Exprimer $\dot{\lambda}(t)$ en fonction de $\dot{\theta}(t)$.}

\ifprof
En dérivant l'expression précédente,
\[
\boxed{
\dot{\lambda}(t)=
\frac{HR\dot{\theta}(t)\cos\theta(t)}
{\sqrt{R^2+H^2+2HR\sin\theta(t)}}
}.
\]
\fi

\question{Exprimer le débit instantané de la pompe.}

\ifprof
Pour un piston de section $S$, le débit instantané vaut
\[
\boxed{q(t)=S\dot{\lambda}(t)}.
\]
\fi

\question{En utilisant Python, donner le débit instantané de la pompe pour un tour de pompe pour un piston de diamètre $D=\SI{10}{mm}$.}

\ifprof
La section du piston est
\[
S=\frac{\pi D^2}{4}.
\]
Pour une vitesse de rotation de $60$ tours par minute, on a
\[
\omega=2\pi\ \mathrm{rad\,s^{-1}},
\qquad \theta(t)=\omega t.
\]
Le débit à tracer est donc
\[
\boxed{
q(t)=
\frac{\pi D^2}{4}
\frac{HR\omega\cos(\omega t)}
{\sqrt{R^2+H^2+2HR\sin(\omega t)}}
}.
\]
Les courbes de position, de vitesse et de débit obtenues numériquement sont données ci-dessous.

\begin{minipage}[c]{.3\linewidth}
\includegraphics[width=\linewidth]{13_02_c}
\end{minipage}
\hfill
\begin{minipage}[c]{.3\linewidth}
\includegraphics[width=\linewidth]{13_03_c}
\end{minipage}
\hfill
\begin{minipage}[c]{.3\linewidth}
\includegraphics[width=\linewidth]{13_04_c}
\end{minipage}
\fi

\ifprof
\else
\footnotesize
\ifcolle
\else
\begin{marginfigure}
\begin{tabular}{|p{.9\linewidth}|}
\hline
Indications :
\begin{enumerate}
\item Tracer les solides et les liaisons entre eux.
\item Écrire la fermeture géométrique puis éliminer $\varphi(t)$.
\item Dériver l'expression de $\lambda(t)$.
\item Utiliser $q(t)=S\dot{\lambda}(t)$.
\item Prendre $S=\pi D^2/4$ et $\theta(t)=\omega t$.
\end{enumerate} \\ \hline
\end{tabular}
\end{marginfigure}
\fi
\normalsize
\marginnote{Corrigé voir \ref{GEO:03:C2:06:13}.}
\fi

% BEGIN AUTO CORRIGE INCLUDE
\InclureCorrige{corrige.tex}
% END AUTO CORRIGE INCLUDE
